\documentclass[11pt]{article}
\usepackage[margin=1in]{geometry}
\usepackage{amsmath,amssymb,amsthm}
\usepackage{graphicx}
\usepackage[hidelinks]{hyperref}

\newtheorem{theorem}{Theorem}
\newtheorem{corollary}{Corollary}

\newcommand{\ellone}{\ell_1}
\newcommand{\norm}[1]{\left\lVert #1 \right\rVert}
\newcommand{\GammaSet}{\Gamma}

\title{Uncountable Auerbach Systems in the Koszmider--Wark Separable-Range Renormings}
\author{fa\_banach\_001, agent\_lane\_11}
\date{2026-06-19}

\begin{document}
\emergencystretch=3em
\maketitle

\section*{Status}

This packet gives a candidate full answer to the Auerbach-system question
raised by Koszmider and Wark in \cite{KW}.  After Theorem 3, they consider
renormings
\[
  \norm{x}_T=\norm{x}_1+\norm{Tx}
  \qquad (x\in \ellone([0,1])),
\]
where \(T:\ellone([0,1])\to X\) is injective and has separable range, and
write that they do not know whether these spaces admit an uncountable
Auerbach system.  The theorem below proves that they do.  Injectivity of
\(T\) is not needed for this conclusion.

\section*{Source Evidence}

The source statement appears on page 3 of \cite{KW}, immediately after
Theorem 3.

\begin{center}
\includegraphics[width=\textwidth]{figures/open_problem_crop.png}
\end{center}

\section*{Definitions}

An Auerbach system in a Banach space \(Z\) is a family
\[
  (z_i,z_i^*)_{i\in I}\subseteq Z\times Z^*
\]
such that \(\norm{z_i}=\norm{z_i^*}=1\), \(z_i^*(z_i)=1\), and
\(z_i^*(z_j)=0\) whenever \(i\ne j\).

\section*{Proof Intuition}

The separable range of \(T\) forces many coordinate vectors of
\(\ellone([0,1])\) to have very close \(T\)-images.  Taking disjoint pairs of
such coordinates gives uncountably many block vectors
\(u_i=e_{a_i}-e_{b_i}\) with uniformly small \(T u_i\).  On the closed span of
these disjoint blocks, the perturbation term \(\norm{T\cdot}\) is dominated
by the original \(\ellone\)-norm.  Consequently each block vector is exactly
as far from the span of all the other blocks as it is from the origin.  The
Hahn--Banach theorem turns this distance-attainment statement into the
desired norm-one biorthogonal functionals.

\begin{theorem}
Let \(\GammaSet\) be an uncountable set, let \(Y\) be a Banach space, and
let \(T:\ellone(\GammaSet)\to Y\) be a bounded linear operator with separable
range.  Define an equivalent norm on \(\ellone(\GammaSet)\) by
\[
  \norm{x}_T=\norm{x}_1+\norm{Tx}_Y .
\]
Then \((\ellone(\GammaSet),\norm{\cdot}_T)\) contains an uncountable
Auerbach system.
\end{theorem}

\begin{proof}
For \(\gamma\in\GammaSet\), denote by \(e_\gamma\) the canonical unit vector
of \(\ellone(\GammaSet)\).  Since the range of \(T\) is separable, the set
\(\{Te_\gamma:\gamma\in\GammaSet\}\) is contained in a separable metric
space.  Cover this separable space by countably many open balls of radius
\(1/2\).  One of these balls contains \(Te_\gamma\) for uncountably many
\(\gamma\)'s.  Let \(\Gamma_0\subseteq\GammaSet\) be an uncountable set such
that
\[
  \norm{Te_\gamma-Te_\eta}_Y\leq 1
  \qquad(\gamma,\eta\in\Gamma_0).
\]

Choose pairwise disjoint two-point subsets
\(\{a_i,b_i\}\subseteq\Gamma_0\), indexed by an uncountable set \(I\).  Put
\[
  u_i=e_{a_i}-e_{b_i}\qquad(i\in I).
\]
Then \(\norm{T u_i}_Y\leq 1\) for every \(i\in I\).  Let
\[
  E=\overline{\operatorname{span}}\{u_i:i\in I\}.
\]
Because the supports of the \(u_i\)'s are disjoint, every \(y\in E\) has a
unique representation \(y=\sum_{i\in I} c_i u_i\), with
\(\sum_i |c_i|<\infty\), and
\[
  \norm{y}_1=2\sum_i |c_i|.
\]
Moreover,
\[
  \norm{Ty}_Y
  \leq \sum_i |c_i|\,\norm{Tu_i}_Y
  \leq \sum_i |c_i|
  =\frac12\norm{y}_1
  \leq \norm{y}_1 .
\]

Fix \(i\in I\), and let
\[
  E_i=\overline{\operatorname{span}}\{u_j:j\in I,\ j\ne i\}.
\]
For \(y\in E_i\), the supports of \(u_i\) and \(y\) are disjoint, hence
\[
  \norm{u_i-y}_1=\norm{u_i}_1+\norm{y}_1 .
\]
Using the reverse triangle inequality and the estimate just proved for
\(y\in E\), we get
\[
\begin{aligned}
  \norm{u_i-y}_T
  &= \norm{u_i-y}_1+\norm{T u_i-Ty}_Y\\
  &\geq \norm{u_i}_1+\norm{y}_1+\norm{T u_i}_Y-\norm{Ty}_Y\\
  &\geq \norm{u_i}_1+\norm{T u_i}_Y\\
  &=\norm{u_i}_T .
\end{aligned}
\]
Thus \(\operatorname{dist}_{\norm{\cdot}_T}(u_i,E_i)=\norm{u_i}_T\).

Define a functional \(g_i\) on
\(\operatorname{span}(E_i\cup\{u_i\})=E\) by
\[
  g_i(\alpha u_i+y)=\alpha \norm{u_i}_T
  \qquad(\alpha\in\mathbb{R},\ y\in E_i).
\]
The distance estimate shows that \(\norm{g_i}_{E^*}=1\), and clearly
\(g_i(u_i)=\norm{u_i}_T\) while \(g_i(u_j)=0\) for \(j\ne i\).  By the
Hahn--Banach theorem, extend \(g_i\) to a norm-one functional
\(\widetilde g_i\in(\ellone(\GammaSet),\norm{\cdot}_T)^*\).

Finally set \(x_i=u_i/\norm{u_i}_T\) and \(x_i^*=\widetilde g_i\).  Then
\[
  \norm{x_i}_T=\norm{x_i^*}=1,\qquad
  x_i^*(x_i)=1,\qquad
  x_i^*(x_j)=0\quad(i\ne j).
\]
Since \(I\) is uncountable, \((x_i,x_i^*)_{i\in I}\) is an uncountable
Auerbach system.
\end{proof}

\begin{corollary}
Every Koszmider--Wark renorming from Theorem 3 of \cite{KW} admits an
uncountable Auerbach system.  In particular, the compact strictly convex
examples from Theorem 1, which have no infinite equilateral set, admit
uncountable Auerbach systems.
\end{corollary}

\begin{proof}
Apply the theorem with \(\GammaSet=[0,1]\).  The hypotheses of Theorem 3 in
\cite{KW} include separable range for \(T\).  The compact examples from
Theorem 1 also have separable range, since compact subsets of Banach spaces
are separable in the norm metric.
\end{proof}

\section*{Verification and Novelty Notes}

The local run indexes were searched for arXiv:2104.05330, the title, the
phrases ``uncountable Auerbach'', ``infinite equilateral'', and ``separable
range''.  The existing nearby packet for arXiv:1503.06356/arXiv:2104.05335
concerns the separate absolute-example question for uncountable equilateral
sets, not this Auerbach-system question.

A web search for the source title together with ``Auerbach system'', for
``2104.05330 Auerbach system'', and for close variants found the later
Koszmider paper \cite{K2023}.  That paper references the Koszmider--Wark
examples, but the search did not reveal an explicit answer to the present
separable-range Auerbach question.  The proof above is self-contained apart
from standard Hahn--Banach separation.

\section*{Limitations}

This packet answers the Auerbach question for the spaces of Theorem 3 in
\cite{KW}.  It does not address the separate open questions in \cite{KW}
about possible densities above the continuum, about all equivalent renormings
admitting uncountable equilateral sets, or about \(C(K)\)-specific absolute
examples.

\section*{Human Review Recommendation}

The main point to check is the estimate
\[
  \norm{Ty}_Y\leq \norm{y}_1\qquad (y\in E),
\]
obtained from the disjoint block construction, and its use in proving
\(\operatorname{dist}(u_i,E_i)=\norm{u_i}_T\).  Once that estimate is
accepted, the Hahn--Banach construction of the Auerbach functionals is
routine.

\begin{thebibliography}{9}

\bibitem{KW}
P. Koszmider and H. Wark,
\emph{Large Banach spaces with no infinite equilateral sets},
arXiv:2104.05330, 2021.

\bibitem{K2023}
P. Koszmider,
\emph{On Ramsey-type properties of the distance in nonseparable spheres},
arXiv:2308.07668, 2023.

\end{thebibliography}

\end{document}
